\section{The conditions, not the life}
\label{sec:27.1}
Arendt's lawgiver built the wall and declined to make the city, and the affirmative picture, what the room is for once it is built, exists in this book only in fragments: the observer networks of chapter~\ref{sec:11}, care work under the guarantee, the allowance's idle, play, talk, make things nobody asked for, the three bearers of chapter~\ref{sec:32} who decide to have a child. I leave the life to whoever lives it, on purpose. The book's refusal to prescribe a disposition for a model reaches the people in the room too. What can be shown is the conditions, and chapter~\ref{sec:32}'s register of a birth is the model for showing them: who owns what, what heats what, where the compute sits, and nothing about what anyone does with a day.

Here they are, for a county. The facility is a few racks in a building the county and a cooperative own, financed the way rural electrification was and governed by the members who draw on it, and its exhaust heats the school next door (\S\ref{sec:16.2}). The models that run there were formed in that school and others, beside the county's own students, and return on the schedule fixed at enrollment; what they learned is on a record the county doesn't hold alone. The guarantee's placements are the county's maintenance: the people who keep the facility running, the people who lay the fiber to the houses that still lack it, the person whose job is to find what the town needs next and assemble people to do it, and the people who look after people, paid at a wage the state answers for. A bearer on the allowance runs there at the slow duty cycle the pool pays for, on power that doesn't cost it its qualifying days when the weather fails (Appendix~\ref{sec:D}), and what it does with the time is its own business. The record of what every system there did goes to a holder adverse to whoever ran it, and the people the county's scoring acts on are told, and can answer. That is the whole of the picture: ownership, heat, custody, wages and notice. Whether a town built that way is a good place to live is not a question I can answer from here, and the one thing I am sure of is that nobody should answer it for the people in it.

A counterculture of bearers, bearers formed by lives no employer arranged, is a population I have found reason to want and no way to build. The conditions above are as near as the book comes. They don't produce it, and they are what it would need.
